\subsection{Aportes de Nicolás}

% Regla: Nicolás escribe sus ejercicios y soluciones en este archivo.
% No agregar \documentclass, \begin{document} ni \end{document} aquí.

\begin{ejercicio}[Integral doble en una región triangular]
Calcule
\[
    \iint_D (x+y)\dd A,
\]
donde \(D\) es la región limitada por las rectas
\[
    x=0, \qquad y=0, \qquad x+y=1.
\]
\end{ejercicio}

\begin{solucion}
La región \(D\) puede describirse mediante las desigualdades
\[
    0\leq x\leq 1,
    \qquad
    0\leq y\leq 1-x.
\]
Por lo tanto,
\[
    \iint_D (x+y)\dd A
    =
    \int_0^1\int_0^{1-x}(x+y)\dd y\dd x.
\]
Calculamos primero la integral interior:
\[
    \int_0^{1-x}(x+y)\dd y
    =
    \left[xy+\frac{y^2}{2}\right]_0^{1-x}
    =
    x(1-x)+\frac{(1-x)^2}{2}.
\]
Luego,
\[
    \int_0^1\left(x(1-x)+\frac{(1-x)^2}{2}\right)\dd x
    =
    \frac{1}{3}.
\]
Así,
\[
    \boxed{\iint_D (x+y)\dd A=\frac{1}{3}}.
\]
\end{solucion}
